\chapter{Cleaning and repairing a structure}
\label{ch:cleaning}

A downloaded PDB file is a record of an experiment, not a ready simulation
system. It usually contains crystal waters, salts from the crystallization
buffer, sometimes extra copies of the protein, and pieces that the experiment
could not see. This chapter shows how to keep only what you want to simulate
(PDB Filter), how to find and repair missing parts (Structure Fixer), and how to
change residues (Amino Acid Mutator).

\begin{glance}
\glanceitem{What it does}{Removes unwanted chains, waters, ions and ligands;
  reports missing residues, missing atoms and alternate conformations; adds
  terminal caps; records and applies mutations.}
\glanceitem{You need it when}{Almost always. Filtering is needed for nearly
  every crystal structure. Repair matters when the file has gaps.}
\glanceitem{Before you start}{A loaded structure (Chapter~\ref{ch:structure}).
  For real repairs and standard mutations: a MODELLER license key.}
\glanceitem{Where to find it}{\menupath{4a}{PDB Filter},
  \menupath{4c}{Structure Fixer}, \menupath{4b}{Amino Acid Mutator}. In
  Workflow Mode they are in the ``fix'' stage.}
\glanceitem{You get}{\file{filtered_structure.pdb},
  \file{repaired_structure.pdb}, and optional reports
  (\file{filter_statistics.json}, \file{completeness_results.json}).}
\end{glance}

\section{The order of the steps}
\label{sec:cleaning-order}

Do the steps in this order: filter (\key{4a}), then, if you need them, mutate
(\key{4b}) and repair (\key{4c}), then protonate (Chapter~\ref{ch:protonation}).
Each tool finds the result of the previous one by itself. For example, the
Structure Fixer prints \prompt{Using filtered structure} when a filtered file
exists, and the Protonation State Analyzer later prints
\prompt{Using repaired structure: \ldots{}/repaired\_structure.pdb}.

\begin{tip}
Read the \prompt{Using \ldots} line at the top of every tool. It tells you which
file the tool is working on. Section~\ref{sec:structure-switching} explains why
this matters.
\end{tip}

Inside a tool's own menu, a filled dot \texttt{●} marks a step you have
already done in this session, a tick \texttt{✓} a step you can do now, and an
open circle \texttt{○} with a \texttt{⚠} line a step that is blocked. The main
menu never shows a tool as done: after filtering, \key{4a} still shows a plain
tick. Main-menu codes such as
\key{1a} do not work inside a tool's menu; type \key{m} first to return to the
main menu.

\section{Step by step: keeping what you need with the PDB Filter}
\label{sec:cleaning-filter}

The example is HIV-1 protease with the inhibitor MK1, \pdb{1HSG}. The file holds
two protein chains (the protease is a dimer), one inhibitor, and 127 crystal
waters. The goal here is the protein dimer alone.

\begin{steps}
  \item Load \pdb{1HSG} (Chapter~\ref{ch:structure}). In the main menu, type
        \key{4a}, then \key{1} (Filter PDB structure).
  \item The first time, ProPrep downloads a dictionary of chemical
        components. Then it measures how the chains touch each
        other:
\begin{screen}
Using RCSB PDB PDB file: .../hsg/1HSG.pdb
...
Chain Interface Analysis
Calculating buried surface area between chains. This may take a moment...
Chain Interface Map
  (Buried Surface
    Area in Ų)
┏━━━┳━━━━━━┳━━━━━━┓
┃   ┃  A   ┃  B   ┃
┡━━━╇━━━━━━╇━━━━━━┩
│ A │  ■   │ 3610 │
│ B │ 3610 │  ■   │
└───┴──────┴──────┘
...
Connectivity (chain ↔ neighbor [BSA, Ų]):
  A ↔ B (3610)
\end{screen}
        The \term{buried surface area} (BSA) is the surface that two chains
        hide from water by touching each other. A large value (over
        1000~Ų; ProPrep calls over 2000~Ų ``very large'') means the chains
        form a real complex. A small value often means a contact made only by
        the crystal. Use it to decide which chains belong together.
  \item ProPrep asks
        \prompt{Detect redox-active sites before filtering? [y/n] (y)}. Answer
        \key{y} if your protein has metal centers or cofactors such as hemes;
        the filter then protects them from accidental deletion
        (Section~\ref{sec:cleaning-protect}). For \pdb{1HSG}, type \key{n}.
  \item ProPrep lists the chains and asks whether to open the 3D viewer
        (default \key{n}). Then choose the chains, by their numbers in the
        list or \typed{all}:
\begin{screen}
Available Chains
1. Chain A - 143 residues - Interfaces with: B (3610Ų)
2. Chain B - 183 residues - Interfaces with: A (3610Ų)

Would you like to view the structure in the interactive 3D viewer? [y/n] (n):
§ans«n»
Select chains (comma-separated indices, or 'all') (all): §ans«all»
\end{screen}
  \item A \prompt{Chain Composition} table follows for each chain. It lists
        every \term{component}, that is, every kind of material in the chain:
        \prompt{Protein}, \prompt{Water}, and each ligand or ion under its
        chemical name. In \pdb{1HSG}, chain B holds the inhibitor, shown by the
        start of its long chemical name
        \prompt{;N-[2(R)-HYDROXY-1(S)-INDANYL]\ldots} and its code
        \prompt{MK1}. Answer \key{y} at
        \prompt{Do you want to proceed with filtering?}.
  \item Now ProPrep asks about each component of each chain in turn. Type
        \key{r} to keep it, \key{d} to remove it, or \key{s} to keep only some
        of its residues:
\begin{screen}
Chain A - Protein
[r] Retain entire component
[s] Select specific residues
[d] Discard entire component
Choose option (r): §ans«r»

Chain A - Water
...
Choose option (r): §ans«d»
...
Chain B - ;N-[2(R)-HYDROXY-1(S)-INDANYL]-5-[(2(S)-TERTIARY
...
Choose option (r): §ans«d»
\end{screen}
        For \pdb{1HSG}: chain A protein \key{r}, chain A water \key{d}, chain B
        protein \key{r}, the inhibitor \key{d}, chain B water \key{d}.
  \item ProPrep summarizes your choices. Check them, then type \key{y}:
\begin{screen}
╭───── Filter Selections ─────╮
│ Model: 0                    │
│                             │
│ Chain A:                    │
│   - Amino_acid: 99 residues │
│ Chain B:                    │
│   - Amino_acid: 99 residues │
│                             │
╰─────────────────────────────╯
Do you want to apply these filters? [y/n] (y): §ans«y»
Filtered structure automatically saved to: filtered_structure.pdb
Structure filtering completed successfully.
\end{screen}
\end{steps}

The new file \file{filtered_structure.pdb} contains the two protein chains
(1514 atoms) and nothing else. It holds only coordinates: the header lines of
the original file (title, REMARKs, sequence records) are not copied. ProPrep
still knows them from the original download.

\begin{tip}
\textbf{What to keep.} A good default for a first simulation:
\begin{itemize}
  \item \textbf{Keep} the protein chains that form the working molecule (see
        the interface table and Section~\ref{sec:structure-asu}).
  \item \textbf{Keep} cofactors and metal ions that the protein needs to
        function or to hold its shape: hemes, iron--sulfur clusters, flavins,
        structural or catalytic zinc, calcium. They usually need extra
        parameters (Chapters~\ref{ch:redox} and~\ref{ch:forcefields}).
  \item \textbf{Usually remove} crystal waters. The simulation surrounds the
        protein with fresh water anyway. Exception: a few buried waters that
        are part of the structure (Section~\ref{sec:cleaning-waters}).
  \item \textbf{Usually remove} ions and small molecules from the
        crystallization buffer, such as chloride, sulfate or glycerol. The
        simulation adds its own salt later (Chapter~\ref{ch:topology}).
  \item \textbf{Keep a ligand} only if you want to study it; it will need
        parameters (Chapter~\ref{ch:forcefields}).
\end{itemize}
\end{tip}

\subsection{Ions: an example with insulin}

In insulin, \pdb{1ZNI}, the ions appear as their own components:
\begin{screen}
Chain B Composition
┏━━━━━━━━━━━━━━━━┳━━━━━━━━━┳━━━━━━━┳━━━━━━━┓
┃                ┃ Residue ┃       ┃ Redox ┃
┃ Component Type ┃  Name   ┃ Count ┃ Sites ┃
┡━━━━━━━━━━━━━━━━╇━━━━━━━━━╇━━━━━━━╇━━━━━━━┩
│ Protein        │ [Total] │  30   │       │
...
│ Water          │ [Total] │  20   │       │
│ CHLORIDE ION   │ [Total] │   3   │       │
│                │   CL    │   3   │       │
│ ZINC ION       │ [Total] │   2   │       │
│                │   ZN    │   2   │       │
└────────────────┴─────────┴───────┴───────┘
\end{screen}
We kept chains A and B (\typed{1,2} at \prompt{Select chains}), removed the
chloride ions with \key{d}, and kept one of the two zinc ions with \key{s}:
\begin{screen}
Chain B - ZINC ION
...
Choose option (r): §ans«s»
             Residues in Chain B - ZINC ION
┏━━━━━━━━┳━━━━━━━━━━━━━┳━━━━━━━━━━━━━━┳━━━━━━━━━━━━━━━━┓
┃ Select ┃ Record Type ┃ Residue Name ┃ Residue Number ┃
┡━━━━━━━━╇━━━━━━━━━━━━━╇━━━━━━━━━━━━━━╇━━━━━━━━━━━━━━━━┩
│ 1      │ HETATM      │ ZN           │ 31             │
│ 2      │ HETATM      │ ZN           │ 32             │
└────────┴─────────────┴──────────────┴────────────────┘
Select residues (comma-separated indices, 'all', or 'none') (all): §ans«1»
\end{screen}
At this question you type the numbers in the \prompt{Select} column, not the
residue numbers.

\begin{watchout}
At \prompt{Select residues}, type only single numbers separated by commas, such
as \typed{1,2,3}, or \typed{all} or \typed{none}. In version 1.23.0 a range such
as \typed{1-50}, or a number larger than the list, stops the filter with an
error message:
\begin{screen}
Error executing PDB Filter.filter_pdb_structure: invalid literal for int()
with base 10: '1-50'
\end{screen}
(or \prompt{list index out of range}). ProPrep returns to the PDB Filter menu,
and \emph{all} answers given so far in this filter run are lost. Start the
filter again.
\end{watchout}

\subsection{Keeping only some waters}
\label{sec:cleaning-waters}

Some waters sit deep inside a protein and are part of its structure. To keep
those and drop the rest, answer \key{s} at a \prompt{Water} component. ProPrep
offers eight ways to analyze the waters:
\begin{screen}
Select Analysis Metrics to Display:
   1. Distance to nearest metal ion
   2. Hydrogen bond analysis
   3. B-factor/thermal information
   4. Water burial: solvent-accessible area and enclosure
   5. Interface proximity
   6. Multi-radius atom-count profiling
   7. Directional atom-count analysis (8 sectors)
   8. Water network analysis (clusters and chains)
   9. Return with no waters selected and proceed with component filtering
Select metrics to display (comma-separated or single choice) (6): §ans«4»
\end{screen}
Type \emph{one} number; despite the wording, \typed{2,4} is refused
(\prompt{Received '2,4', which is not one of the available options}). Choice
\key{4} (water burial) is the most direct: after an optional change of the
probe size, it prints a table of all waters, most buried first. For chain A of
\pdb{1HSG}, one of 44 waters, \prompt{324}, was \prompt{Enclosed} (sealed inside
the protein, next to Arg87); all others were \prompt{Exposed}. A long
explanation of each column is printed with the table.

Then type \key{3} (\prompt{Proceed with water selection}) and give the waters
to keep \emph{by their residue numbers} from the table:
\begin{screen}
Available: 305 to 575 (44 waters total)

Select waters by residue numbers (all): §ans«324,305,307»

Keep 3 waters, discard 41 waters.
Keeping: HOH305, HOH307, HOH324
Confirm selection? [y/n] (y): §ans«y»
\end{screen}
Every number must be a water. A range is refused if any number in it is not a
water (\prompt{Invalid residue numbers: [306, 308, 310]}); you are then asked
again. Choice \key{9} at the metrics menu keeps no waters.

\subsection{Protecting cofactors: redox-site detection}
\label{sec:cleaning-protect}

If you answer \key{y} to \prompt{Detect redox-active sites before filtering?},
ProPrep first runs the Redox Site Detector (Chapter~\ref{ch:redox}) and marks
what it finds. A component that holds such a site then gets a different
question:
\begin{screen}
⚠  This component contains redox site residues:
   • Site site_1 (): MK1902
[r] Retain entire component
[c] Keep redox site residues
[s] Select specific residues
[o] Override redox protection and discard (⚠ breaks redox sites)
Choose option (r):
\end{screen}
\key{c} keeps only the residues that belong to the site. \key{o} removes them
anyway. Within one session the question is asked only once; later filter runs
reuse the sites found.

\subsection{Checking and saving what was removed}

In the PDB Filter menu, \key{2} (Show filtering status) shows the original and
filtered residue counts and what was removed, for example
\prompt{Hetero groups: MK1, HOH}. \key{3} (Export filter selection) saves the
same information as \file{filter_statistics.json}.

\begin{watchout}
\begin{itemize}
  \item Answering \key{n} at \prompt{Do you want to apply these filters?}
        returns to the menu without a message and without saving.
  \item Filtering again overwrites \file{filtered_structure.pdb} without
        asking. When the workspace holds several structures, the filter first
        asks which one to use, in a table (for example Oriented,
        Protonation-Updated, Filtered, RCSB PDB). Choose the original download
        if you want to start again.
  \item The filter does not remove \term{alternate locations} (see
        Section~\ref{sec:cleaning-fixer}); the filtered \pdb{1ZNI} still had
        10 atoms with two positions.
\end{itemize}
\end{watchout}

\section{Step by step: finding and repairing gaps with the Structure Fixer}
\label{sec:cleaning-fixer}

Experiments often miss parts of a protein. Flexible loops and chain ends may be
invisible, so whole residues are absent (\term{missing residues}). Some side
chains are only partly seen (\term{missing atoms}). And some atoms are seen in
two or more positions at once, each with a share of the time called its
\term{occupancy} (\term{alternate locations}, or altlocs). A simulation needs
exactly one complete position for every atom. The Structure Fixer finds all
three problems and repairs what it can.

\begin{needsoftware}
Real repairs need \term{MODELLER}, a separate program for building protein
models, with a free academic license key. Without a valid key, ProPrep prints
this message at every start:
\begin{screen}
MODELLER not found. Structure repair (missing loops/residues) will be
unavailable.
  To install: conda install -c salilab modeller
  A free academic license is required:
https://salilab.org/modeller/registration.html
  After registering, set your license key:
    export KEY_MODELLER='your_license_key_here'
\end{screen}
The message appears also when MODELLER is installed and only the key is missing
or wrong, as on our test computer.

\textbf{Works without the key:} all the detection and reports; adding
\term{terminal caps} (ACE and NME); applying non-standard mutations from
\key{4b}, but only when you also add at least one cap.

\textbf{Does not work without the key:} building missing residues and loops;
rebuilding missing side-chain atoms; standard mutations (for example Lys to
Arg); and choosing one alternate location. The menu says so with
\prompt{MODELLER off: capping only}.

To add a key, ProPrep reads the environment variable \texttt{KEY\_MODELLER}, or
the first line of the file \file{~/.proprep/modeller_key} (not tried for this
manual). Appendix~\ref{app:software} and Chapter~\ref{ch:installation} have
the details.
\end{needsoftware}

The example is \pdb{2RH1}, the $\beta_2$-adrenergic receptor. It is a
\term{chimera}: part of the receptor was replaced by the protein T4 lysozyme to
help it crystallize. Its file lists missing residues in REMARK 465 lines.

\begin{steps}
  \item Filter the structure first (Section~\ref{sec:cleaning-filter}). Then
        type \key{4c}. Options \key{2}--\key{4} stay blocked until you have
        analyzed the structure. Type \key{1}
        (\prompt{Analyze structure completeness}).
  \item ProPrep asks three questions, one per kind of problem. Press
        \key{Enter} three times to accept the defaults, which read the
        experimenters' own records in the file:
\begin{screen}
Missing Residues Detection:
1. REMARK 465 records
2. SEQRES comparison
3. FASTA comparison
4. Sequence gaps
5. Skip
Choose (1): §ans«1»

Missing Atoms Detection:
1. REMARK 470 records
2. Template comparison
3. Skip
Choose (1): §ans«1»

Alternate Locations Detection:
1. AltLoc identifiers
2. Occupancy values
3. Skip
Choose (1): §ans«1»
\end{screen}
        \term{REMARK 465} lists residues that were in the sample but not seen;
        \term{REMARK 470} lists side-chain atoms that were not seen. Option
        \key{2} for missing residues compares the atoms with the
        \term{SEQRES} records, the full sequence of the sample written in the
        file; for \pdb{2RH1} it found the same 58 residues.
  \item Read the summary. For \pdb{2RH1}:
\begin{screen}
┃ Issue Type                     ┃ Method            ┃ Count ┃ Chains ┃
┡━━━━━━━━━━━━━━━━━━━━━━━━━━━━━━━━╇━━━━━━━━━━━━━━━━━━━╇━━━━━━━╇━━━━━━━━┩
│ Missing Residues               │ remark_465        │ 58    │ A      │
│ Missing Atoms                  │ remark_470        │ 1     │ A      │
│ Alternate Locations            │ altloc_identifier │ 0     │ -      │
│ Pending Standard Mutations     │ user_specified    │ 0     │ -      │
│ Pending Non-Standard Mutations │ user_specified    │ 0     │ -      │
\end{screen}
        Below it, the sequence of each chain is printed with colors: green
        for missing residues, cyan for residues with missing atoms, yellow for
        alternate locations, magenta and orange for pending mutations, and
        \prompt{-} for numbers that are not in the sequence at all.
  \item Answer \key{y} at \prompt{Show detailed reports?} to see each missing
        residue and atom. For \pdb{2RH1}, the 58 missing residues are at the
        two ends of the chain (positions $-6$ to 28 and 343 to 365). The one
        residue with missing atoms is Asp29.
\end{steps}

\begin{watchout}
Points that can mislead you when reading the results in version 1.23.0:
\begin{itemize}
  \item In a chimera or engineered construct, the sequence view shows long
        runs of \prompt{-} where the numbering jumps (in \pdb{2RH1}, between
        230 and 263, and up to 1002 where the T4 lysozyme numbering starts).
        These are \emph{not} missing residues.
  \item The REMARK 470 detail table shows the chain letter in the
        \prompt{Atom} column (\prompt{ASP 29 | A}). The file itself says the
        missing atoms of Asp29 are CG, OD1 and OD2.
  \item \prompt{Template comparison} (option \key{2} for missing atoms) reports
        a missing \prompt{OXT} atom on \emph{every} residue: 75 for ubiquitin,
        about 440 for \pdb{2RH1}. The \term{OXT} atom (the second oxygen of the
        last residue's carboxyl group) exists only at the end of a chain, so
        these reports are false. Use REMARK 470 for files from the RCSB.
        Template comparison also downloads templates first, which took 1--2
        minutes.
  \item \prompt{FASTA comparison} returned 0 missing residues for \pdb{2RH1}
        and then said \prompt{No structural issues found!}, although REMARK 465
        lists 58. Do not use it for entries with one chain.
  \item \prompt{Sequence gaps} counts every jump in numbering as missing
        residues: 691 for \pdb{2RH1}. It is not useful for constructs with
        non-continuous numbering.
\end{itemize}
\end{watchout}

\subsection{Applying repairs and caps}
\label{sec:cleaning-caps}

The two ends of a protein chain carry charges: a positive amino group at the
start (the \term{N-terminus}) and a negative carboxyl group at the end (the
\term{C-terminus}). When your structure starts or ends where the real protein
continues (because residues are missing, or because you cut out a domain),
these artificial charges can distort the simulation. A \term{cap} is a small
neutral group that closes the chain end instead: \term{ACE} (an acetyl group)
before the first residue, and \term{NME} (an N-methyl amide) after the last.

\begin{steps}
  \item In the Structure Fixer menu, type \key{2}
        (\prompt{Apply repairs/mutations/caps to structure}). Without MODELLER
        you go straight to capping:
\begin{screen}
MODELLER not available — gap filling and mutations are disabled.
Terminal capping (ACE/NME) is still available; install MODELLER for full
repairs.
...
═══ Terminal Capping (ACE/NME) ═══
ACE caps the N-terminus, NME caps the C-terminus. Cap atoms get placeholder
geometry that tLEaP finalizes during topology generation.

Chain A — N-term: ASP29, C-term: TYR1161
     n  N-terminus only — ACE before ASP29
     c  C-terminus only — NME after TYR1161
  both  Both termini
     s  Skip this chain (default)
Cap which terminus of chain A? (s): §ans«n»
\end{screen}
  \item Type \key{n}, \key{c} or \typed{both} for each chain. Pressing
        \key{Enter} (\key{s}) adds no cap; if no chain gets a cap, ProPrep says
        \prompt{No terminal caps selected.} and writes nothing.
\begin{screen}
Adding ACE/NME caps...
  Plan: ACE before A:29 (MODELLER numbering, capping N-terminal gap 29-29)
  Inserted ACE at A:28 (before 29)
✓ Caps added successfully

✓ Structure repair completed!
Output: repaired_structure.pdb
\end{screen}
\end{steps}

The file \file{repaired_structure.pdb} keeps all ligands and waters that were in
the input. Without MODELLER, the missing loops and termini are \emph{not} built:
they remain as breaks in the chain. The missing side-chain atoms of Asp29 are
not rebuilt either; tLEaP normally adds missing side-chain atoms when it builds
the system (Chapter~\ref{ch:topology}).

\begin{watchout}
\textbf{The repaired file is renumbered from 1.} In the repaired \pdb{2RH1},
the cap is residue 1, Asp29 became residue 2, and the last residue, Tyr1161,
became 506. In ubiquitin with an ACE cap, every number moves up by one (Lys48
becomes 49). All later steps use the new numbers: the protonation tables
(Chapter~\ref{ch:protonation}), residue overrides, and orientation by residue.
Keep a note of the shift when you compare with the literature.
\end{watchout}

\begin{watchout}
Without MODELLER, alternate locations are not resolved. On \pdb{3NIR}
(crambin at 0.48~Å), the analysis found 32 residues with alternates, for
example \prompt{TYR 29 \quad A, B, C \quad A:0.42 B:0.30 C:0.28}. ProPrep has a
screen for choosing one conformation, but in version 1.23.0 it is reached only
when MODELLER works. The repaired file still contained 313 atoms with A and B
positions. Most crystal structures at ordinary resolution have few or no
altlocs; check the summary line \prompt{Alternate Locations}.
\end{watchout}

\subsection{Reports}

\key{3} (View structural issues report) prints the detailed tables again.
\key{4} (Export results to JSON) saves them as \file{completeness_results.json}
(including any false OXT reports).

\section{Changing residues with the Amino Acid Mutator}
\label{sec:cleaning-mutator}

Use the mutator to simulate a mutant, for example ubiquitin with Lys48 changed
to arginine. You can also put in a modified residue, for example a
phosphoserine or the infrared probe cyanophenylalanine. The mutator only
\emph{records} the changes. The Structure Fixer applies them later.

There are two kinds of mutation:
\begin{itemize}
  \item A \textbf{standard mutation} changes one of the 20 natural amino acids
        into another. MODELLER builds the new side chain.
  \item A \textbf{non-standard mutation} changes a residue into anything else,
        named by its three-letter code (for example \prompt{SEP}
        phosphoserine, \prompt{TPO} phosphothreonine, \prompt{PTR}
        phosphotyrosine, \prompt{MLY} dimethyllysine, \prompt{CNF}
        cyanophenylalanine). You choose which atoms of the old residue to keep;
        tLEaP builds the rest later. The new residue must exist in your force
        field (Chapter~\ref{ch:forcefields}).
\end{itemize}

\begin{needsoftware}
Standard mutations need MODELLER with a valid license key. Without it, ProPrep
records them but never applies them, and the Structure Fixer does not tell you
so when it writes the repaired file. The suggestion line in the Structure Fixer
menu warns beforehand: \prompt{Found 2 pending mutations, but MODELLER is
required to apply them}. Non-standard mutations work without MODELLER, but only
if you add at least one terminal cap at the same time.
\end{needsoftware}

\begin{steps}
  \item Type \key{4b}. First type \key{1}
        (\prompt{Display protein sequences with numbering}). The mutator does
        not let you enter mutations before you have seen the sequence.
\begin{screen}
Chain A:
                 10        20        30        40        50        60
   1-60  MQIFVKTLTGKTITLEVEPSDTIENVKAKIQDKEGIPPDQQRLIFAGKQLEDGRTLSDYN

                 10
  61-76  IQKESTLHLVLRLRGG

Found 1 protein chain(s) with 76 total residues
\end{screen}
        The ruler above each line counts from that line's first residue: the
        \prompt{10} above line 61--76 marks residue 70.
  \item Type \key{2} (\prompt{Specify amino acid mutations}). Enter each
        mutation as \textsf{CHAIN:POSITION:FROM->TO}, one per line:
\begin{screen}
Enter mutation: §ans«A:48:LYS->ARG»
✓ Added standard mutation: Chain A, Position 48, LYS->ARG

Enter mutation: §ans«A:63:LYS->ALA»
✓ Added standard mutation: Chain A, Position 63, LYS->ALA

Enter mutation: §ans«A:4:PHE->CNF»

⚠ Detected non-standard residue target: CNF
Non-standard mutations require atom selection (atoms to keep from original
residue)
...
Preset Options:
  backbone  - Keep backbone atoms only (N, CA, C, O)
  bb+cb     - Keep backbone + CB (beta carbon)
  all       - Keep all atoms (for testing)
  custom    - Manually select specific atoms
  cancel    - Cancel this mutation

Select preset or 'custom': §ans«backbone»
✓ Added non-standard mutation: Chain A, Position 4, PHE->CNF
  Keeping 4 atoms: N, CA, C, O
\end{screen}
        For a non-standard residue, ProPrep lists the atoms of the old residue.
        \typed{bb+cb} is a good choice when the new residue keeps the beta
        carbon (as in phosphoserine). With \typed{custom} you type atom names,
        for example \typed{N,CA,C,O,CB,OG}.
  \item Type \typed{done}. ProPrep shows two tables, standard and
        non-standard, and says
        \prompt{Automatically stored mutations in workspace!}.
  \item Go to \key{4c}, analyze (\key{1}), then apply (\key{2}). The pending
        mutations appear in the summary. Add at least one cap:
\begin{screen}
Cap which terminus of chain A? (s): §ans«both»
...
═══ Non-Standard Mutations ═══

Applying 1 non-standard mutation(s)...

→ Mutating A:4 PHE → CNF
  Mapped residue: A:4 → A:5
    Removed 7 atoms
    Renamed residue to CNF
  ✓ Mutation applied successfully
    Kept 4 atoms: N, CA, C, O
...
Output: repaired_structure.pdb
\end{screen}
        Note the renumbering: with the ACE cap, Phe4 became CNF~5.
\end{steps}

Without MODELLER, the two standard mutations (K48R and K63A) were \emph{not}
applied: the repaired file still has lysine at those positions, and \key{4b}
option \key{3} still lists them as waiting.

The mutator checks what you type and asks again after a mistake:

\begin{tabularx}{\linewidth}{@{}lL@{}}
\toprule
\textbf{You type} & \textbf{ProPrep says} \\
\midrule
\typed{K48R} & \prompt{Invalid format. Use CHAIN:POSITION:FROM->TO} \\
\typed{A:48:ALA->GLY} & \prompt{Position 48 is LYS, not ALA} \\
\typed{A:200:ALA->GLY} & \prompt{Position 200 not found in chain A} \\
\typed{B:48:LYS->ARG} & \prompt{Chain B not found in structure} \\
\typed{A:48:LYS->GLU} (again) & \prompt{Mutation already specified for chain A position 48} \\
\bottomrule
\end{tabularx}

\medskip
Other commands at \prompt{Enter mutation}: \typed{help} shows more examples and
common non-standard codes; \typed{clear} deletes \emph{all} pending mutations at
once, without asking for confirmation.

\begin{watchout}
\begin{itemize}
  \item If you add no cap, a pending non-standard mutation is silently not
        applied: ProPrep says \prompt{No terminal caps selected.} and writes
        nothing, and the summary even says \prompt{No pending mutations}.
  \item After a non-standard mutation, \key{4b} option \key{1} leaves the new
        residue out of the one-letter sequence (75 letters instead of 76).
\end{itemize}
\end{watchout}

\section{Options at a glance}

\begin{longtable}{@{}>{\raggedright\arraybackslash}p{0.38\linewidth}>{\raggedright\arraybackslash}p{0.57\linewidth}@{}}
\toprule
\textbf{Question} & \textbf{Answers} \\
\midrule
\endhead
\multicolumn{2}{@{}l}{\textsf{\bfseries 4a PDB Filter}} \\
Detect redox-active sites? & \key{y} (default)
  protects metal sites and cofactors; \key{n} skips \\
\prompt{Select chains} & \typed{1,2} or \typed{all} (default) \\
Each component & \key{r} keep (default), \key{s} choose residues,
  \key{d} remove; for a redox site also \key{c} keep site residues,
  \key{o} remove anyway \\
\prompt{Select residues} & \typed{1,3,5} (table numbers), \typed{all},
  \typed{none}; no ranges \\
Water metrics & one number 1--8; \key{9} keeps no waters \\
\prompt{Select waters by residue numbers} & residue numbers and ranges,
  \typed{all}, \typed{none} \\
Apply these filters? & \key{y} (default); \key{n}
  discards silently \\
\midrule
\multicolumn{2}{@{}l}{\textsf{\bfseries 4c Structure Fixer}} \\
Missing residues & \key{1} REMARK 465 (default, recommended), \key{2} SEQRES,
  \key{3} FASTA, \key{4} numbering gaps, \key{5} skip \\
Missing atoms & \key{1} REMARK 470 (default, recommended), \key{2} template
  comparison (false OXT reports), \key{3} skip \\
Alternate locations & \key{1} altloc letters (default), \key{2} occupancy
  values, \key{3} skip \\
\prompt{Cap which terminus} & \key{n}, \key{c}, \typed{both}, \key{s} (default,
  none) \\
\midrule
\multicolumn{2}{@{}l}{\textsf{\bfseries 4b Amino Acid Mutator}} \\
\prompt{Enter mutation} & \textsf{A:48:LYS->ARG}, \typed{done}, \typed{clear},
  \typed{help} \\
\prompt{Select preset or 'custom'} & \typed{backbone}, \typed{bb+cb},
  \typed{all}, \typed{custom}, \typed{cancel} \\
\bottomrule
\end{longtable}

\section{If something goes wrong}
\label{sec:cleaning-problems}

\begin{description}[style=nextline]
  \item[The filter stops with ``invalid literal for int()''] You typed a range
        at \prompt{Select residues}. Start again and type single numbers
        separated by commas.
  \item[``list index out of range'' in the filter] A number was larger than
        the list. Start again.
  \item[``Received '2,4', which is not one of the available options''] The
        water-metrics question accepts one number only.
  \item[``Received '1a', which is not one of the available options''] You
        typed a main-menu code inside a tool. Type \key{m} first.
  \item[The fixer says ``capping only''] MODELLER or its license key is
        missing. See the box in Section~\ref{sec:cleaning-fixer}.
  \item[The repaired file has no mutation] Standard mutations need MODELLER.
        Non-standard ones need at least one cap.
  \item[Residue numbers do not match the paper] The repaired file was
        renumbered from 1 (Section~\ref{sec:cleaning-caps}).
  \item[Hundreds of ``OXT'' atoms reported missing] These are false reports
        of the template comparison. Use REMARK 470.
  \item[The repaired file still has A and B positions] Alternate locations
        are not resolved without MODELLER.
\end{description}

See also Chapters~\ref{ch:troubleshooting} and~\ref{ch:knownissues}.
